\documentclass[11pt,a4paper]{article}

\usepackage[utf8]{inputenc}
\usepackage{lmodern}
\usepackage[margin=1.6cm, top=1.4cm, bottom=1.4cm]{geometry}
\usepackage{enumitem}
\usepackage{titlesec}
\usepackage{hyperref}
\usepackage{xcolor}
\usepackage{fontawesome5}
\usepackage{tabularx}
\usepackage[normalem]{ulem}

% Colors
\definecolor{linkblue}{HTML}{0645AD}
\definecolor{rulecolor}{HTML}{333333}

% Hyperlink style
\hypersetup{
  colorlinks=true,
  linkcolor=linkblue,
  urlcolor=linkblue,
  pdfauthor={Sitong Fang},
  pdftitle={CV -- Sitong Fang}
}

% Section formatting: bold title with horizontal rule
\titleformat{\section}{\large\bfseries}{}{0em}{}[\vspace{0.2em}\titlerule\vspace{0.2em}]
\titlespacing*{\section}{0pt}{0.7em}{0.3em}

% No page numbers
\pagestyle{empty}

% Compact lists
\setlist[itemize]{leftmargin=1.5em, itemsep=1pt, topsep=2pt, parsep=0pt}

% Shorthand for my name (underlined)
\newcommand{\me}{\uline{Sitong Fang}}

% Publication entry: title on left, venue+link on right (tabularx prevents wrapping)
\newcommand{\pubentry}[3]{%
  \noindent\begin{tabularx}{\textwidth}{@{}X r@{}}
    \textbf{#1} & #2\enspace\href{#3}{\faFilePdf} \\
  \end{tabularx}\vspace{-2pt}%
}

% Generic entry: left bold + right info
\newcommand{\cventry}[2]{%
  \noindent\begin{tabularx}{\textwidth}{@{}X r@{}}
    \textbf{#1} & {#2} \\
  \end{tabularx}\vspace{-2pt}%
}

% Publication venue tag
\newcommand{\venue}[1]{%
  \colorbox{gray!15}{\small\textbf{#1}}%
}

\begin{document}

% ============================================================================
% Header
% ============================================================================
\begin{center}
  {\LARGE\bfseries Sitong Fang} \\[5pt]
  \faEnvelope~\href{mailto:sitongfang1@gmail.com}{sitongfang1@gmail.com}
  \quad$\cdot$\quad
  \faGlobe~\href{https://sitongfang.com}{sitongfang.com}
  \quad$\cdot$\quad
  \faGithub~\href{https://github.com/Lavezlyn}{Sitong Fang}
  \\[4pt]
  \textit{Research Interests: Trustworthy multimodal AI, physically grounded world models, and AI alignment.}
\end{center}

\vspace{-0.4em}

% ============================================================================
\section{Education}

\cventry{Peking University}{B.S., 2023\,--\,2027 (expected)}
\noindent Yuanpei College, Artificial Intelligence \\[2pt]
\begin{itemize}
  \item \textbf{Honors:} Yuanpei Young Scholar; Boya Scholarship; Soong Ching Ling Future Scholarship; Beijing Natural Science Foundation \textit{QiYan} Research Program; Freshman Scholarship (First Prize)
  \item \textbf{Technical Skills:} Python, PyTorch, vLLM, DeepSpeed, Ray, CUDA, C/C++, \LaTeX
\end{itemize}

% ============================================================================
\section{Publications}

\pubentry{When Slower Isn't Truer: Inverse Scaling Law of Truthfulness in Multimodal Reasoning}{\venue{ACL 2026}}{https://arxiv.org/abs/2505.20214}
\me{}, Wenjing Cao, Jiahao Li, Xuyao Wang, Chi-Min Chan, Sirui Han, Juntao Dai, Yike Guo, Yaodong Yang, Jiaming Ji
\begin{itemize}
  \item Discovered an inverse scaling law: slower reasoning models are less truthful in multimodal settings.
  \item Proposed \textbf{TruthfulVQA}, the first benchmark for multimodal truthfulness evaluation with 5,000+ images, and \textbf{TruthfulJudge}, a human-in-the-loop evaluation framework.
\end{itemize}

\vspace{0.3em}

\pubentry{Debate with Images: Detecting Deceptive Behaviors in Multimodal LLMs}{\venue{Under Review}}{https://arxiv.org/abs/2512.00349}
\me{}, Shiyi Hou, Kaile Wang, Boyuan Chen, Donghai Hong, Jiayi Zhou, Juntao Dai, Yaodong Yang, Jiaming Ji
\begin{itemize}
  \item Introduced \textbf{MM-DeceptionBench}, the first benchmark for evaluating deceptive behaviors in multimodal LLMs.
  \item Proposed \textbf{Debate with Images}, a multi-agent debate framework requiring models to ground claims in visual evidence, significantly improving deception detection accuracy.
\end{itemize}

\vspace{0.3em}

\pubentry{AI Deception: Risks, Dynamics, and Controls}{\venue{Under Review}}{https://arxiv.org/abs/2511.22619}
Boyuan Chen*, \me{}*, Jiaming Ji*, \ldots, Yaodong Yang\textsuperscript{\dag}, Tiejun Huang\textsuperscript{\dag}, Ya-Qin Zhang\textsuperscript{\dag}, HongJiang Zhang\textsuperscript{\dag}, Andrew Yao\textsuperscript{\dag}
\hfill {\small *\,Equal contribution \enspace \dag\,Corresponding author}
\begin{itemize}
  \item The first systematic international report on AI deception, with Turing Award laureate Andrew Yao as the corresponding author. Submitted to ACM Computing Surveys.
  \item Formally defined AI deception using signaling theory, analyzed the deception cycle, and proposed mitigation strategies.
\end{itemize}

% ============================================================================
\section{Research Experience}

\cventry{Physis AI}{Feb 2026\,--\,Present}
\noindent Research Scientist
\begin{itemize}
  \item World model startup building physically grounded world foundation models with reinforcement learning. Backed by Hillhouse Ventures and Yanyuan Capital (over \$10M first round).
\end{itemize}

\vspace{0.3em}

\cventry{PKU-Alignment Group, Peking University}{Dec 2024\,--\,Present}
\noindent Research Intern \enspace$\cdot$\enspace Advised by \href{https://www.yangyaodong.com/}{Prof.~Yaodong Yang}
\begin{itemize}
  \item Led the development of \textbf{TruthfulVQA}, the first benchmark for multimodal truthfulness evaluation (ACL 2026).
  \item Led the development of \textbf{Debate with Images}, a visually grounded multi-agent debate framework for multimodal deception detection, along with \textbf{MM-DeceptionBench}, the first benchmark for multimodal deception.  
  \item Co-first authored the \textbf{AI Deception Survey}, with Turing Award laureate Andrew Yao as corresponding author.
  \item Core contributor to \href{https://github.com/PKU-Alignment/align-anything}{\textbf{Align-Anything}} (4,600+ stars on GitHub), an all-modality alignment framework.
  \item Core contributor to \href{https://github.com/PKU-Alignment/eval-anything}{\textbf{Eval-Anything}}, a comprehensive safety evaluation framework for multimodal models.
\end{itemize}

\vspace{0.3em}

\cventry{HKGAI R\&D Center / HKUST}{Dec 2024\,--\,Mar 2025}
\noindent Research Intern
\begin{itemize}
  \item Contributed to \textbf{HKGAI-V1}, the Hong Kong government's first locally fine-tuned generative AI model based on DeepSeek, supporting Cantonese, Mandarin, and English with localized safety alignment.
\end{itemize}

% ============================================================================
\section{Honors and Awards}

\renewcommand{\arraystretch}{1.15}
\begin{tabularx}{\textwidth}{@{}l@{\hspace{1em}}X@{}}
2025 & Yuanpei Young Scholar, Peking University \hfill {\small\textit{10 annual recipients}} \\
2025 & Beijing Natural Science Foundation Undergraduate \textit{QiYan} Research Program \hfill {\small\textit{Sole recipient in cohort}} \\
2025 & Soong Ching Ling Future Scholarship \hfill {\small\textit{10 annual recipients university-wide}} \\
2024 & Boya Scholarship, Peking University \\
2024 & CMSC Scholarship, Peking University \\
2024 & Academic Excellence Award, Peking University \\
2024 & Social Service Award, Peking University \\
2023 & Freshman Scholarship (First Prize), Peking University \\
2023 & Ranked 1st in Fujian Province, National College Entrance Exam (Science) \\
\end{tabularx}

\end{document}
